\documentclass[10pt,a4paper]{amsart}
\usepackage[margin=19mm]{geometry}
\usepackage{amsmath,amssymb,mathtools}
\usepackage[scr=rsfs,cal=euler]{mathalfa}
\usepackage{xfrac}
\usepackage[unicode,hidelinks,bookmarks=false]{hyperref}
\newcommand{\ie}{i.e.\ }
\newcommand{\cf}{cf.\ }
\newcommand{\resp}{respectively\ }
\newcommand{\Subsection}{\subsection}
\providecommand{\nobreakdash}{\nobreak\hbox{-}}
\setlength{\emergencystretch}{2em}
\multlinegap0pt
\usepackage{amssymb}
\newtheorem{thm}{Theorem}
\newcommand{\C}{\mathbb {C}}     % Complex  number field
\title{Reducibility of symmetry on $q$-Painlev\'e equations}
\author{Chihiro Sato, Kouichi Takemura, Ohka Yamashita}
\date{Source: arXiv:2609.02927v1 (CC BY 4.0)}
\begin{document}
\maketitle

\noindent\emph{Source and licence.} Reducibility of symmetry on $q$-Painlev\'e equations. Version arXiv:2609.02927v1 (CC BY 4.0), distributed by arXiv under the Creative Commons Attribution 4.0 International licence, http://creativecommons.org/licenses/by/4.0/, as recorded in the arXiv licence metadata for this exact version and shown on the abstract page https://arxiv.org/abs/2609.02927v1. Full licence notice: \texttt{licenses/arxiv-2609.02927-CC-BY-4.0.txt}.\par\smallskip

\noindent\textit{Editorial scope.} Counted unit: the article's thm thm:E6compati (source0.tex lines 514--517). The article's complete original proof of it is delivered with it (source0.tex lines 518--572, one \texttt{proof} environment of 55 lines). No mathematics is written, repaired or re-derived: the statement and the proof are the article's own lines, in the article's own notation, and no retained line is substituted or re-flowed.  Retained uncounted as the source-local context the proof needs: lines 156--171; lines 505--510.  Nothing was omitted from any retained span, so the package carries no omission record.  No retained line contains a citation, so the wrapper carries no bibliography and no bibliography entry.  No retained line contains a figure environment, an image-inclusion command or drawing code, so no image is delivered and the package declares no asset dependency.  Editorial formatting only, in addition: an AMS \texttt{amsart} preamble that restates the article's own declarations and macros used by the retained lines, namely the environments \texttt{thm} on separate counters, together with the standard packages those lines need.  The retained environments print in this document as Theorem 1 in order of appearance, whereas the article prints the same items with the numbering of its own full document.  The complete native source of the article as distributed in the arXiv e-print for this exact version is bundled unchanged as \texttt{sources/209251/source0.tex}.
\begin{align}
s_0 : & \: \nu_7 \leftrightarrow \nu_8 , \quad s_1 : \: \nu_5 \leftrightarrow \nu_6  , \quad s_3 : \: \nu_1 \leftrightarrow \nu_2 , \quad s_4 : \: \nu_2 \leftrightarrow \nu_3 , \quad s_5 : \: \nu_3 \leftrightarrow \nu_4 , \label{eq:AffineWeylActionE6} \\
 s_2 : & \: \nu_1 \rightarrow\frac{\kappa_2}{\nu_6}, \;
 \nu_6 \rightarrow \frac{\kappa_2}{\nu_1}, \;
 \kappa_1 \rightarrow \frac{\kappa_1\kappa_2}{\nu_1\nu_6}, \;
 f \rightarrow f\, \frac{\kappa_2 ( \nu_1 g-1)  }{ (\nu _1 \nu_6 - \kappa_2  )fg + \nu _1 \kappa_2 g - \nu _1 \nu_6 } ,\nonumber \\
 s_6 : & \: \nu_1 \rightarrow\frac{\kappa_1}{\nu_7}, \; 
 \nu_7 \rightarrow \frac{\kappa_1}{\nu_1}, \;
 \kappa_2 \rightarrow \frac{\kappa_1\kappa_2}{\nu_1\nu_7}, \;
  g \rightarrow g\, \frac{\nu _7 ( \nu_1 - f )  }{ \kappa_1  -  \nu _7 f + ( \nu _1 \nu_7 - \kappa_1 )fg } , \nonumber \\
  \pi _1 : & \: q \rightarrow 1/q, \: \nu_1 \rightarrow \nu_2 / \kappa_2 , \:  \nu_2 \rightarrow \nu_1 / \kappa_2 ,  \: \nu_3 \rightarrow 1/\nu_6,  \: \nu_4 \rightarrow 1/\nu_5, \: \nu_5 \rightarrow 1/\nu_4,  \nonumber \\
& \: \nu_6 \rightarrow 1/\nu_3, \: \nu_7 \rightarrow 1/\nu_7, \: \nu_8 \rightarrow 1/\nu_8, \: \kappa_1 \rightarrow \nu_1 \nu_2 /(\kappa_1 \kappa_2 ) , \: \kappa_2 \rightarrow 1/\kappa_2 ,\nonumber \\
& \: f \rightarrow \frac{\nu _1 \nu_2 ( 1- fg )}{ \kappa_2 \{ \nu _1 \nu_2 g + f- (\nu _1 + \nu _2 ) fg \} } , \: g \rightarrow \kappa_2 g , \nonumber \\
  \pi _2 : & \: q \rightarrow 1/q, \: \nu_1 \rightarrow 1/\nu_1, \:  \nu_2 \rightarrow 1/\nu_2,  \: \nu_3 \rightarrow 1/\nu_3,  \: \nu_4 \rightarrow 1/\nu_4, \: \nu_5 \rightarrow 1/\nu_8, \nonumber \\
& \: \nu_6 \rightarrow 1/\nu_7, \: \nu_7 \rightarrow 1/\nu_6, \: \nu_8 \rightarrow 1/\nu_5, \: \kappa_1 \rightarrow 1/\kappa_2 , \: \kappa_2 \rightarrow 1/\kappa_1 , \: f \leftrightarrow  g . \nonumber 
\end{align}

\begin{align}
&  (\nu '_1 , \nu '_2 , \nu '_3 , \nu '_4 ,\nu '_5 ,\nu '_6 , \nu '_7 ,\nu '_8 ; \kappa '_1 ,\kappa '_2 ; f' ,g' ) \sim  (\nu_1 , \nu_2 , \nu_3 , \nu_4 ,\nu_5 ,\nu_6 , \nu_7 ,\nu_8 ;\kappa_1 ,\kappa_2 ; f ,g ) \label{eq:equivE6} \\
& \Leftrightarrow \mbox{There exist } a, b, c \in \C (\nu_1 , \nu_2 , \cdots ,\nu_8 ,\kappa_1 ,\kappa_2 , f ,g ) \setminus \{ 0 \} \mbox{ such that }  \nonumber \\
& \qquad (\nu '_1 , \nu '_2 , \nu '_3 , \nu '_4 ,\nu '_5 ,\nu '_6 , \nu '_7 ,\nu '_8 ; \kappa '_1 ,\kappa '_2 ; f' ,g' )  \nonumber \\
& \qquad \qquad = (a \nu_1 , a \nu_2 , a \nu_3 , a \nu_4 ,b \nu_5 , b \nu_6 , c \nu_7 , c \nu_8 ; ac \kappa_1 , ab \kappa_2 ; af ,g/a) .  \nonumber
\end{align}

\begin{thm} \label{thm:E6compati}
Let $(\nu '_1 , \nu '_2 , \nu '_3 , \nu '_4 ,\nu '_5 ,\nu '_6 , \nu '_7 ,\nu '_8 ; \kappa '_1 ,\kappa '_2 ; f' ,g' ) $ be a tuple of the parameters which is equivalent to $ (\nu_1 , \nu_2 , \nu_3 , \nu_4 ,\nu_5 ,\nu_6 , \nu_7 ,\nu_8 ;\kappa_1 ,\kappa_2 ; f ,g )$.
The actions of the generators of  $\widetilde{W}(E^{(1)}_6)$ in equation (\ref{eq:AffineWeylActionE6}) to $(\nu '_1 , \nu '_2 , \nu '_3 , \nu '_4 ,\nu '_5 ,\nu '_6 , \nu '_7 ,\nu '_8 ; \kappa '_1 ,\kappa '_2 ; f' ,g' ) $ are written as the same form of the actions to $(\nu _1 , \nu _2 , \nu _3 , \nu _4 ,\nu _5 ,\nu _6 , \nu _7 ,\nu _8 ; \kappa _1 ,\kappa _2 ; f ,g ) $ up to the equivalence.
\end{thm}

\begin{proof}
By the assumption, there exist $ a, b, c \in \C (\nu_1 , \nu_2 , \cdots ,\nu_8 ,\kappa_1 ,\kappa_2 , f ,g ) \setminus \{ 0 \} $ such that 
\begin{align}
&  (\nu '_1 , \nu '_2 , \nu '_3 , \nu '_4 ,\nu '_5 ,\nu '_6 , \nu '_7 ,\nu '_8 ; \kappa '_1 ,\kappa '_2 ; f' ,g' )  \label{eq:equivE6'} \\
& = (a \nu_1 , a \nu_2 , a \nu_3 , a \nu_4 ,b \nu_5 , b \nu_6 , c \nu_7 , c \nu_8 ; ac \kappa_1 , ab \kappa_2 ; af ,g/a). \nonumber
\end{align}
We confirm the theorem on the action of $\pi _1$.
Recall that
\begin{align}
& (\pi _1 ( \nu _1) , \pi _1 (\nu _2) , \pi _1 ( \nu _3) , \pi _1 ( \nu _4) , \pi _1 (\nu _5) , \pi _1 (\nu _6) , \pi _1 ( \nu _7) , \pi _1 (\nu _8) ;  \pi _1 (\kappa _1) , \pi _1 (\kappa _2) ;  \pi _1 (f ) , \pi _1 (g) ) \\
& = ( \frac{\nu  _2 }{\kappa _2}  , \frac{\nu  _1 }{\kappa _2}   , \frac{1}{ \nu _6} , \frac{1}{\nu  _5}  , \frac{1}{\nu  _4}  ,  \frac{1}{\nu  _3}  ,  \frac{1}{\nu _7} , \frac{1}{\nu _8} ; \frac{\nu _1 \nu _2 }{ \kappa_1\kappa _2}  , \frac{1}{ \kappa _2} ; \frac{\nu _1 \nu_2 ( 1- fg )}{ \kappa_2 \{ \nu _1 \nu_2 g + f- (\nu _1 + \nu _2 ) fg \} } , \kappa_2 g  ). \nonumber 
\end{align}
We calculate the action of $\pi _1$ to $\nu '_1 , \nu '_2 , \dots  , f' ,g' $ which satisfy equation (\ref{eq:equivE6'}).
We have
\begin{align}
&  \pi _1 ( \nu '_5 ) = \pi _1 ( b ) \pi _1 (\nu _5 )  =  \frac{\pi _1 ( b ) }{\nu _4} = \frac{\pi _1 ( b ) a }{ \nu '_4 } , \; \pi _1 ( f ') =  \pi _1 ( a ) \pi _1 (f )  \\
&  \quad = \frac{ \pi _1 ( a ) \nu _1 \nu_2 ( 1- fg )}{ \kappa_2 \{ \nu _1 \nu_2 g + f- (\nu _1 + \nu _2 ) fg \} }  = \frac{ \pi _1 ( a ) b \nu '_1 \nu '_2 ( 1- f'g' )}{ \kappa ' _2 \{ \nu '_1 \nu'_2 g' + f' - (\nu '_1 + \nu '_2 ) f' g' \} }  , \nonumber
\end{align}
and so on.
It follows from a direct calculation that
\begin{align}
%& (\pi _1 ( \nu '_1) , \pi _1 (\nu '_2) , \pi _1 ( \nu '_3) , \pi _1 ( \nu '_4) , \pi _1 (\nu '_5) , \pi _1 (\nu '_6) , \pi _1 ( \nu '_7) , \pi _1 (\nu '_8) ; \pi _1 (\kappa '_1) , \pi _1 (\kappa '_2) ;  \pi _1 (f' ) , \pi _1 (g') ) \\
& (\pi _1 ( \nu '_1) , \pi _1 (\nu '_2) , \pi _1 ( \nu '_3) , \pi _1 ( \nu '_4) , \dots , \pi _1 (\nu '_8) ; \pi _1 (\kappa '_1) , \pi _1 (\kappa '_2) ;  \pi _1 (f' ) , \pi _1 (g') ) \\
& = ( \pi _1 ( a ) b \frac{ \nu '_2 }{ \kappa ' _2} , \pi _1 ( a ) b \frac{ \nu '_1 }{\kappa ' _2} ,  \frac{ \pi _1 ( a ) b  }{ \nu '_6 }  , \frac{ \pi _1 ( a ) b  }{ \nu '_5 } , \frac{  \pi _1 ( b ) a }{ \nu '_4 } , \frac{ \pi _1 ( b ) a }{ \nu '_3} , \frac{ \pi _1 (c ) c }{ \nu '_7} , \frac{ \pi _1 ( c ) c }{ \nu '_8} ; \nonumber \\
& \qquad \pi _1 ( a c) b c  \frac{  \nu '_1 \nu '_2 }{ \kappa '_1\kappa '_2}  , \frac{ \pi _1 ( ab ) ab }{\kappa '_2 }  ; \pi _1 ( a ) b  \frac{\nu '_1 \nu '_2 ( 1- f'g' )}{ \kappa ' _2 \{ \nu '_1 \nu'_2 g' + f' - (\nu '_1 + \nu '_2 ) f' g' \} } , \frac{\kappa '_2 g' }{\pi _1 ( a ) b }  ) , \nonumber 
\end{align} 
which is equivalent to
\begin{align}
& ( \frac{ \nu '_2 }{ \kappa ' _2} , \frac{ \nu '_1 }{\kappa ' _2} ,  \frac{1 }{ \nu '_6 }  , \frac{ 1 }{ \nu '_5 } , \frac{1 }{ \nu '_4 } , \frac{ 1}{ \nu '_3} , \frac{ 1 }{ \nu '_7} , \frac{ 1 }{ \nu '_8} ;  \frac{  \nu '_1 \nu '_2 }{ \kappa '_1\kappa '_2}  , \frac{ 1 }{\kappa '_2 }  ;   \frac{\nu '_1 \nu '_2 ( 1- f'g' )}{ \kappa ' _2 \{ \nu '_1 \nu'_2 g' + f' - (\nu '_1 + \nu '_2 ) f' g' \} } , \kappa '_2 g'  ) ,
\end{align} 
where the parameters $\pi _1 ( a ) b  ,  \pi _1 ( b ) a ,  \pi _1 ( c ) c$ correspond to the parameters $a,b,c $ in equation (\ref{eq:equivE6}).
Therefore, the action of $\pi _1 $ to $(\nu '_1 , \nu '_2 ,$ $ \nu '_3 , \nu '_4 ,\nu '_5 ,\nu '_6 , \nu '_7 ,\nu '_8 ; \kappa '_1 ,\kappa '_2 ; f' ,g' ) $ is written as the same form of the action to $(\nu _1 , \nu _2 , \nu _3 ,$ $ \nu _4 ,\nu _5 ,\nu _6 , \nu _7 ,\nu _8 ; \kappa _1 ,\kappa _2 ; f ,g ) $ up to the equivalence.
It is shown similarly that the action of the other generators of $\widetilde{W}(E^{(1)}_{6})$ to $(\nu '_1 , \nu '_2 ,$ $ \nu '_3 , \nu '_4 ,\nu '_5 ,\nu '_6 , \nu '_7 ,\nu '_8 ; \kappa '_1 ,\kappa '_2 ; f' ,g' ) $ is written as the same form of the action to $(\nu _1 , \nu _2 , \nu _3 ,$ $ \nu _4 ,\nu _5 ,\nu _6 , \nu _7 ,\nu _8 ; \kappa _1 ,\kappa _2 ; f ,g ) $ up to the equivalence.
We note crucial equations on the actions of some generators of $\widetilde{W}(E^{(1)}_{6})$.
\begin{align}
& (\pi _2 ( \nu _1) , \pi _2 (\nu _2) , \pi _2 ( \nu _3) , \pi _2 ( \nu _4) , \dots , \pi _2 (\nu _8) ;  \pi _2 (\kappa _1) , \pi _2 (\kappa _2) ;  \pi _2 (f ) , \pi _2 (g) ) \\
& = (\frac{1}{\nu_1}, \frac{1}{\nu_2}, \frac{1}{\nu_3}, \frac{1}{\nu_4}, \frac{1}{\nu_8} , \frac{1}{\nu_7}, \frac{1}{\nu_6}, \frac{1}{\nu_5} ; \frac{1}{\kappa_2} , \frac{1}{\kappa_1}  ; g , f ),  \nonumber \\
& (\pi _2 ( \nu '_1) , \pi _2 (\nu '_2) , \pi _2 ( \nu '_3) , \pi _2 ( \nu '_4) , \dots , \pi _2 (\nu '_8) ;  \pi _2 (\kappa '_1) , \pi _2 (\kappa '_2) ;  \pi _2 (f' ) , \pi _2 (g') ) \nonumber \\
& =  (\frac{\pi _2 ( a ) a }{\nu ' _1}, \frac{\pi _2 ( a ) a }{\nu  '_2}, \frac{\pi _2 ( a ) a }{\nu  '_3}, \frac{\pi _2 ( a ) a }{\nu  '_4}, \frac{\pi _2 ( b ) c }{\nu  '_8} , \frac{\pi _2 ( b ) c }{\nu  '_7}, \frac{\pi _2 ( c ) b }{\nu  '_6}, \frac{\pi _2 ( c ) b }{\nu  '_5} ; \nonumber \\
& \qquad   \frac{\pi _2 ( a c ) a b }{\kappa  '_2} , \frac{\pi _2 ( a b ) a c }{\kappa  '_1 } ; \pi _2 ( a ) a g ' ,  \frac{f ' }{\pi _2 ( a ) a}) . \nonumber
\end{align} 
\begin{align}
& (s _2 ( \nu _1) , s _2 (\nu _2) , s _2 ( \nu _3) , s _2 ( \nu _4) , \dots , s _2 (\nu _8) ;  s _2 (\kappa _1) , s _2 (\kappa _2) ;  s _2 (f ) , s _2 (g) ) \\
& =  (\frac{\kappa_2 }{\nu_6 }, \nu_2 , \nu_3 , \nu_4 ,\nu_5 ,\frac{\kappa_2 }{\nu_1} , \nu_7 ,\nu_8 ; \frac{\kappa_1 \kappa_2 }{\nu_1 \nu_6} ,\kappa_2 ;  f\, \frac{\kappa_2 ( \nu_1 g-1)  }{ (\nu _1 \nu_6 - \kappa_2  )fg + \nu _1 \kappa_2 g - \nu _1 \nu_6 } ,g ) , \nonumber \\
& (s _2 ( \nu '_1) , s _2 (\nu '_2) , s _2 ( \nu '_3) , s _2 ( \nu '_4) , \dots , s _2 (\nu '_8) ;  s _2 (\kappa '_1) , s _2 (\kappa '_2) ; s _2 (f' ) , s _2 (g') ) \nonumber \\
& =  (\frac{s _2 (a)}{a} \frac{\kappa '_2 }{\nu '_6} , \frac{s _2 (a)}{a} \nu '_2 , \frac{s _2 (a)}{a} \nu '_3 , \frac{s _2 (a)}{a} \nu '_4 , \frac{s _2 (b)}{b} \nu '_5 , \frac{s _2 (b)}{b} \frac{\kappa '_2 }{\nu '_1} , \frac{s _2 (c)}{c} \nu '_7 , \frac{s _2 (c)}{c} \nu '_8 ; \nonumber \\
& \qquad \frac{s _2 (ac)}{ac} \frac{\kappa '_1 \kappa '_2 }{\nu '_1 \nu '_6} , \frac{s _2 (ab)}{ab} \kappa '_2 ; \frac{s _2 (a)}{a} f' \frac{\kappa '_2 ( \nu '_1 g' -1)}{(\nu ' _1 \nu '_6 - \kappa '_2  ) f' g' + \nu ' _1 \kappa '_2 g' - \nu ' _1 \nu '_6 } , \frac{a}{s _2 (a)} g' ) . \nonumber 
\end{align} 
\begin{align}
& (s _6 ( \nu _1) , s _6 (\nu _2) , s _6 ( \nu _3) , s _6 ( \nu _4) , \dots , s _6 (\nu _8) ;  s _6 (\kappa _1) , s _6 (\kappa _2) ;  s _6 (f ) , s _6 (g) ) \\
& =  (\frac{\kappa_1 }{\nu_7} , \nu_2 , \nu_3 , \nu_4 ,\nu_5 ,\nu_6 , \frac{\kappa_1 }{\nu_1} ,\nu_8 ; \kappa_1, \frac{\kappa_1 \kappa_2 }{\nu_1 \nu _7 } ;  f, \frac{\nu _7 ( \nu_1 - f ) g }{ \kappa_1  -  \nu _7 f + ( \nu _1 \nu_7 - \kappa_1 )fg } ) , \nonumber \\
& (s _6 ( \nu '_1) , s _6 (\nu '_2) , s _6 ( \nu '_3) , s _6 ( \nu '_4) , \dots , s _6 (\nu '_8) ; s _6 (\kappa '_1) , s _6 (\kappa '_2) ; s _6 (f' ) , s _6 (g') ) \nonumber \\& =  (\frac{s _6 (a)}{a} \frac{\kappa '_1 }{\nu '_7} , \frac{s _6 (a)}{a} \nu '_2 , \frac{s _6 (a)}{a} \nu '_3 , \frac{s _6 (a)}{a} \nu '_4 , \frac{s _6 (b)}{b} \nu '_5  , \frac{s _6 (b)}{b} \nu '_6 , \frac{s _6 (c)}{c} \frac{\kappa '_1 }{\nu '_1} , \frac{s _6 (c)}{c} \nu '_8 ; \nonumber \\
& \qquad \frac{s _6 (ac)}{ac} \kappa '_1, \frac{s _6 (ab)}{ab} \frac{\kappa '_1 \kappa '_2 }{\nu '_1 \nu '_7 } ;  \frac{s _6 (a)}{a} f', \frac{a}{s _6 (a)} \frac{\nu '_7 ( \nu '_1 - f' ) g' }{ \kappa '_1  -  \nu '_7 f' + ( \nu '_1 \nu '_7 - \kappa '_1 )f'g' } ) . \nonumber 
\end{align} 
\end{proof}

\end{document}
